% -*- mode: latex; coding: latin-1-unix -*- %

\documentclass[a4paper,11pt]{article}



\usepackage[utf8]{inputenc} % pour les accents
\usepackage[T1]{fontenc}
\usepackage[francais]{babel}
\usepackage{fullpage}
\usepackage{color}
\usepackage{amsmath}
\usepackage{url} 
\usepackage{makeidx}
\usepackage{amsmath}
\usepackage{wasysym}
\usepackage{float}
\usepackage{graphicx}
\usepackage{a4wide}
\usepackage{multirow}
\usepackage{afterpage}
\usepackage{pst-all} % Pour pstricks
\usepackage{moreverb}

\definecolor{gray}{gray}{0.85}

\title{
  \normalsize{\begin{flushright} MASTER 2 - Semestre 2 - 2009 \end{flushright}}
  \vspace{15mm}
  \hrule height 1mm
  \vspace{5mm}
  \Huge{\emph{Projet de fin d'étude\\\textsl{AltaRica et les robots.}}}
  \vspace{5mm}\hrule height 1mm
  \vspace{1cm}
}

\author{
  J\'{e}r\'{e}my BOYER\\
  Dorian CRANSAC\\
  Jean S\'{e}bastien DUBERNET\\
  Fabien KUNTZ\\
  \vspace{2cm}
}

\date{}

\begin{document}
\begin{titlepage}
  \begin{figure}
    \vspace{1cm}
    \psset{unit=1in,linewidth=4pt} %paramtrage des units pour pstricks
    \rput(1,0){\includegraphics[width=4cm]{Bordeaux1}}
    \vspace{15mm}
  \end{figure}
\end{titlepage}

\maketitle

\vspace{4cm}
\begin{center}
Enseignant : Alain GRIFFAULT
\end{center}

\tableofcontents

\newpage
\begin{abstract}

Nous allons vous pr\'{e}senter dans ce rapport notre projet de fin
d'\'{e}tudes, qui consiste en l'application du cours de conceptions
formelles sur la programmation de robots et la g�n�ration automatique
de code. Les robots en question sont des robots NXT de l'entreprise
LEGO. Nous utilisons le langage NXC pour les programmer. Ces robots
sont capables d'int�ragir avec leur environnement gr\^{a}ce \`{a} un
ensemble de capteurs et de moteurs. Nous utilisons Altarica pour
montrer que les diff�rents comportements des robots sont s�res.\newline

Nous avons d'abord du nous familiariser avec les robots NXT,
les possibilités des différents capteurs, moteurs et
constructions. Les possibilités au niveau code ont aussi été
éxaminées afin de savoir quel type de mission est réalisable. Nous
avons ensuite choisi la mission.\newline

Ensuite, nous avons d�velopp� en parall�le le code correspondant � la
mission et le mod�le Altarica. De nombreuses mises-�-jour ont �t�
n�cessaires avant d'obtenir les versions finales. Nous avons alors
tent� de prouver grâce à une spécification $Acheck$ pour chaque
noeud Altarica qu'il est s�r et ne provoque pas de comportements ind�sirables. \newline

Enfin, nous avons d�crit les m�canismes principaux formant les bases
d'une g�n�ration automatique de code NXC � partir d'un mod�le
altarica. Nous n'avons pas pu �crire le parseur correspondant en
raison de la contrainte de temps li�e au projet mais restons convaincu
qu'un tel outil est r�alisable � partir de notre travail.



\end{abstract}





 \section{R\'{e}sum\'{e}}

 intro...

 \section{Mission}

  \subsection{Environnement}
  L'environnement est constitu\'{e} d'un marquage au sol. Ce marquage
  correspond \`{a} un chemin d\'{e}limit\'{e} par des feuilles de couleurs
  diff\'{e}rentes de celle du sol. Un point d'arriv\'{e}e est pr\'{e}d\'{e}fini. On peut
  repr\'{e}senter l'environnement par une grille sur laquelle les
  d\'{e}placements possibles sont verticaux ou horizontaux. Les \'{e}l\'{e}ments de
  marquage sont donc de forme carr\'{e}e et on d\'{e}finit un chemin unique
  repr\'{e}sentant la combinaison de d\'{e}placements que les robots doivent
  \^{e}tre capables de r\'{e}aliser. Les couleurs des cases de la grille sont
  altern\'{e}es de mani\`{e}re \`{a} obtenir un damier. Ce d\'{e}tail permet au robot de
  diff\'{e}rencier une impasse d'un chemin possible lorsqu'il \'{e}value les
  possibilit\'{e}s autour de sa position courante.

  \subsection{Entit\'{e}s}

   \subsubsection{Robot Ma\^{i}tre}
   Le robot ma\^{i}tre est le robot qui poss\`{e}de les capteurs. C'est
   \'{e}galement lui qui d\'{e}cide des actions que les deux robots doivent
   effectuer. Il est capable de donner des ordres au robot esclave
   gr\^{a}ce \`{a} des composants Bluetooth. Deux capteurs sont mont\'{e}s sur
   ce robot : un capteur de lumi\`{e}re \`{a} l'avant, et un capteur
   ultrasonique \`{a} l'arri\`{e}re. Les d\'{e}placements possibles du robot sont
   les suivants : avancer jusqu'\`{a} la d\'{e}tection d'un changement
   d'environnement, tourner \`{a} gauche ou \`{a} droite sur place \`{a}
   90 deg. Le
   capteur de lumi\`{e}re permet de d\'{e}tecter les bordures du marquage. Le
   capteur ultrasonique detecte la pr\'{e}sence du robot esclave \`{a}
   l'arri\`{e}re. Il est particuli\`{e}rement utile apr\`{e}s un changement de
   direction pour arr\^{e}ter le robot esclave \`{a} hauteur du virage. 

   \subsubsection{Robot Esclave}
   Ce robot ne comporte aucun capteur, et est compl\`{e}tement d\'{e}pendant du
   robot ma\^{i}tre, \`{a} savoir qu'il n'effectue une action que lorsqu'il
   re\c coit un message par la communication Bluetooth. Les possibilit\'{e}s de
   d\'{e}placement de ce robot sont les m\^{e}mes.

   \subsubsection{Int\'{e}ractions}
   Lors du d\'{e}placement dans l'environnement et de la d\'{e}couverte
   d'obstacles, le robot ma\^{i}tre fait des choix de d\'{e}placement et
   communique au robot esclave des ordres lui permettant de passer
   \'{e}galement ces obstacles:
   \begin{itemize}
    \item Le robot ma\^{i}tre avance jusqu'\`{a} la d\'{e}tection d'une bordure
    \item Si une bordure est d\'{e}tect\'{e}e, les d\'{e}cisions suivantes sont
	  prises:
    \begin{itemize}
     \item Le robot tourne \`{a} gauche \`{a} 90 deg.
     \item Si une bordure est pr\'{e}sente, il sait que le chemin qu'il doit
	   emprunter est celui de droite.
    \end{itemize}
   \end{itemize}

  \subsection{Objectif}
  Les robots sont cens\'{e}s se d\'{e}placer \`{a} l'int\'{e}rieur du marquage et
  atteindre le point d'arriv\'{e}e. Le robot ma\^{i} doit explorer le terrain
  et permettre au robot esclave de le suivre. 

\section{R\'{e}duction API NXC}
Pour la r\'{e}alisation du  projet nous utiliserons 4 cat\'{e}gories de fonction : celles relatives aux moteurs, aux capteurs , à la connexion Bluetooth et \'{e}ventuellement d'autres fonctions permettant la coh\'{e}sion de l'ensemble.

\subsection{Les moteurs }
\begin{verbatim}
	OnFwd("port","power"); //permet de mettre en marche "avant" un moteur du robot 
	OnRev("port","power"); //sens inverse
	RotateMotor("port","speed","degrees"); //
	RotateMotorEx("port","speed","degrees","turnpct","sync","stop"); 
\end{verbatim}

\subsection{Les capteurs }

\subsubsection{Type et mode des capteurs}
\begin{verbatim}
	SetSensorLight("port"); //Lie le capteur de lumière à un "port"
	SetSensorSound("port"); //Lie le capteur ultrason à un "port"
	SetSensorMode("port","const mode"); //Configure le mode du capteur associé au "port"
	SetInput("port","const field","value"); //Configure le type du capteur associé au "port"
	SetSensorLowspeed("port"); //Configure le capteur en mode I2C
	ClearSensor("port"); //Efface la valeur d'un capteur associé à un "port"

\end{verbatim}

\subsubsection{Information des capteurs}
\begin{verbatim}
	Sensor("n"); //retourne la valeur d'un capteur
	SensorUS("n"); //retourne la valeur du capteur ultrason

\end{verbatim}

\subsection{La connexion Bluetooth }
\begin{verbatim}
	BluetoothStatues("NUMBER"); //Vérifie la connexion
	RemoteStartProgram("connection", "filename");
	RemoteStopProgram("connection");

\end{verbatim}

\subsection{Autres fonctions }
\begin{verbatim}
	Wait("time"); //Temps donné à l'exécution d'une action
	off("port"); //Stop le moteur associé au "port"
	Acquire("mutex"); 
	Release("mutex"); 
	Precedes("task1","task2","taskn"); //Permet l'exécution de taches simultané

\end{verbatim}

\subsection{Les variables }
\begin{verbatim}
	OUT_A; //port A
	OUT_AC; //port AC
	IN_1; //Input 1
	IN_2; //Input 2

\end{verbatim}

\section{Implémentation}

\subsection{Le langage NXC}

Le langage NXC est un langage calqué sur le C d'où il tire d'ailleurs son
sigle, Not eXactly C. Il est ce que l'on peut appeler une surcouche du
langage utilisé par les robots  : le langage NBC (Next Byte Code). Bien
que celui-ci possède une API assez developpée, ça syntaxe s'inspirant
fortement de l'assembleur n'en reste pas moins rude et c'est pour cela
que nous préfèrerons programmer les robots  dans le langage NXC.


\subsection{Algorithme}


L'algorithme utilisé dans le robot maître est assez simple à comprendre
en lui même. La difficulté rencontré vient donc plus de l'approche de la
programmation du robot.
\\
Pour rapidement décrire l'algorithme, nous avons en premier lieu un while nous
permettant de boucler à l'infini. Dans celui-ci le capteur de lumière
vient mettre à jour une variable, current\_val, qui indique la couleur de la prochaine
case rencontrée par le robot.

\begin{verbatim} 
 while(current_val == last_val){
   last_val = current_val ;
   real_val = Sensor(S3);
  if((real_val>=MIN_GRIS) && (real_val<=MAX_GRIS))
    current_val = GRIS;
  else{
    if((real_val>=MIN_BLANC) && (real_val<=MAX_BLANC))
     current_val = BLANC;
     else if((real_val>=MIN_NOIR) && (real_val<=MAX_NOIR))
       current_val = NOIR;
      }
  } 
\end{verbatim}

Le robot maître avance ensuite jusqu'à la prochaine case. En fonction de la
couleur de celle-ci, deux cas se présentent.
Si elle est grise ou blanche, le robot maître s'arrète et ordonne au
robot esclave d'avancer jusqu'à lui.

\begin{verbatim}
if((changed_square == true) && (obstacle_mode == false))
{
 OnFwd(OUT_AC,0);
 RemoteStartProgram(1,"forward.rxe");
 while(SensorUS(S4) > 10);
 RemoteStopProgram(1);
 OnFwd(OUT_AC,30);
}
\end{verbatim}

Si elle est noire, le robot maître effectue l'opération de contournement
de l'obstacle.

\begin{verbatim} 
if((current_val == NOIR)){
 if(obstacle_mode == FALSE){
 TextOut(10,LCD_LINE4,"NOIR 1");
 obstacle_mode = TRUE;
 OnFwd(OUT_AC,0);
 RotateMotorEx(OUT_AC, 30, 220, -100, true, true);
 }else{
    turned_to_right = TRUE;
    TextOut(10,LCD_LINE5,"NOIR 2");
    OnFwd(OUT_AC,0);
    RotateMotorEx(OUT_AC, 30, 400, 100, true, true);
    obstacle_mode = FALSE;
    }
}
\end{verbatim}

Il ordonne ensuite au robot esclave d'avancer et tourner dans le même
sens. 

\begin{verbatim}
if(found_way == TRUE){
  if(turned_to_right == TRUE){
   OnFwd(OUT_AC,30);
   Wait(3000);
   OnFwd(OUT_AC,0);
  }
  RemoteStartProgram(1,"forward.rxe");
  while(SensorUS(S4) > 10);
  RemoteStopProgram(1);
 
  Wait(2000);
 
  if(turned_to_right == FALSE){
   RemoteStartProgram(1,"turnleft.rxe");
   Wait(5000);
   RemoteStopProgram(1);
  }
  else{
     RemoteStartProgram(1,"turnright.rxe");
     Wait(5000);
     RemoteStopProgram(1);
     }
  found_way = FALSE;
  OnFwd(OUT_AC,30);
  }
\end{verbatim}

Les deux robots recommencent ensuite le même cycle de fonctionement.

 \section{Noeuds AltaRica}
 
  \subsection{LightSensor}
  \verbatiminput{../src/altarica/LightSensor.alt}

  \subsection{UltraSonicSensor}
  \verbatiminput{../src/altarica/UltraSonicSensor.alt}

  \subsection{Motor}
  \verbatiminput{../src/altarica/Motor.alt}

  \subsection{Moving}
  \verbatiminput{../src/altarica/Moving.alt}


\section{Génération de code à partir des noeuds altarica élémentaires.}

Etant donné que notre modèle altarica représente les comportements des éléments du système, nous allons ici faire le lien entre chaque comportement adopté par une entité et des portions de code implémentant ces comportements.

\subsection{Motor}

\begin{verbatim}
event
    	back		OnFwd(Out_X, Y); pour le moteur X, et Y choisi arbitrairement entre 1 et 100.
    	stop		OnFwd(Out_X, 0);
    	forward	OnRev(Out_X, Y); pour le moteur X, et Y choisi arbitrairement entre 1 et 100.

\end{verbatim}

La conjonction des évenements des 2 moteurs permettent de déplacer le robot de manière cohérente. Ceci correspond au noeud Moving.

\subsection{Moving}

\begin{verbatim}
event
	stop	OnFwd(Out_AC, 0);
	left	OnFwd(Out_C, 20); et OnRev(Out_A, 20); ou bien RotateMotorEx(OUT_AC, 40, 180, -100, true, true); (plus simple)
	forward	OnFwd(Out_AC, 40);
	right		OnFwd(Out_A, 20); et OnRev(Out_C, 20); ou bien RotateMotorEx(OUT_AC, 40, 180, 100, true, true); (plus simple)
	halfturn	RotateMotorEx(OUT_AC, 40, 360, 100, true, true);
\end{verbatim}

\subsection{LightSensor}

\begin{verbatim}
flow
    value : [0,2];	value = SensorUs(3); 3 étant dans notre cas le numéro du port sur lequel est branché le capteur lumineux.
\end{verbatim}

\subsection{UltraSonicSensor}

\begin{verbatim}
flow
    d : bool;	value = SensorUs(4); 4 étant dans notre cas le numéro du port sur lequel est branché le capteur ultrason.

\end{verbatim}

\subsection{BlueTooth}
% -*- mode: latex; coding: latin-1-unix -*- %
\section{Bilan}

\subsection{Difficultées rencontrées}

\subsubsection{Problématique du damier}

Nous avions initialement prévu que la mission se déroulerait dans un environnement ou le sol serait homogène (exemple: tout blanc). Cependant, un problème de décision est survenu. Il était impossible de différencier la présence de l'absence d'un chemin à gauche ou à droite lors de la rencontre d'un obstacle. Etant donné que le robot reste sur la même case lors de sa rotation, il aurait fallu décider d'une distance ou durée avant la rencontre potentielle d'un nouvel obstacle sur le nouveau chemin, ce qui va à l'encontre  de l'objectif de réactivité et généricité du comportement du robot. Pour éviter de temporiser cette action, nous avons décidé d'utiliser un damier gris et blanc, permettant de décider au premier changement de case sur le nouveau chemin si ce chemin est effectivement utilisable ou si la solution était la direction opposée (rencontre immédiate d'une nouvelle case noire).

\subsubsection{Calibrage du capteur lumineux}

Selon les différentes intensités lumineuses dans la pièce où se trouve le robot, les bornes des 3 seuils utilisés doivent être adaptées. Tout au long du projet, il a fallu ajuster ces valeurs de manière à ce que le robot détecte bien les différentes couleurs de cases et à ce que les prises de décision ne soient pas faussées. Il était parfois difficile de faire la différence entre un bug logique (erreur dans le code) et un bug physique (intéractions avec l'environnement). Nous estimons que les bornes sont actuellement positionnées de manière à rendre la mission possible dans un maximum d'endroits et de luminosités différentes.

\subsubsection{Problématique de seuil intermédiaire du capteur lumineux}

Lors du passage du capteur entre les 2 seuils extrêmes (case blanche à case noire ou case noire à case blanche), le capteur détecte un seuil intermédiaire gris pendant un bref instant, causant un bug dans la prise de décision. Nous avons du mettre en place un procédé de correction de la detection qui renouvelle la lecture du capteur après un très bref instant et change la valeur obtenue si elle est différente de la dernière lecture. Cette temporisation était incontournable.

\subsubsection{Caractère incontrolable des valeurs retournées par les capteurs}

Au début de notre modélisation Altarica des capteurs, nous pensions pouvoir nous en sortir avec une variable d'état $value$ pour représenter les valeurs "lues" par le capteur de lumière et le capteur ultrasonique. Et un évenement $readValue$ devait survenir en fonction de chaque changement de valeur de notre variable d'état. Cela ne convenait pas ensuite dans la vérification, car en réalité, il est impossible de prévoir les valeurs qui vont \^{e}tres lues par ces capteurs. Le robot est réactif à son environnement. Nous avons donc remplacés ces modélisations par une simple variable de flux, qui est bien incontrolable et donc cohérente.

\subsubsection{Possibilité de tirer des transitions non voulues}

Au moment de simuler notre n\oe{}ud Controller, nous avons remarqué qu'il nous était possible de faire tourner le robot maître à droite, alors qu'en réalité, cela ne correspondait pas. Nous avons ensuite compris qu'il s'agissait d'une action non capturée d'un sous-n\oe{}ud de notre Controller. Nous avons donc remarqué que cet événement n'était pas capturé parce que le robot maître ne tournait jamais à droite.

\subsection{Conclusion}

Nous avons apprécié l'approche que nous avons eue pendant ce projet. Nous ne nous sommes pas uniquement concentrés sur des difficultés logicielles, et avons abordés les problèmes techniques découlant d'un système complet (électronique et logiciel). Nous aurions voulu aller au bout de notre démarche mais le temps nous a manqué.

Il nous semble que la réalisation d'un outil de génération de code automatique, reposant sur notre travail, est tout à fait envisageable en utilisant une méthode proche de celle que nous avons proposée.

\end{document}
